\documentclass[../../../include/open-logic-section]{subfiles}

\begin{document}

\olfileid{sth}{ord-arithmetic}{mult}

\olsection{Ping-pingan Ordinal}

Saiki kita ngalih menyang ping-pingan ordinal, kanthi cara kang
padha karo panambahan ordinal. Upamané kita arep ngitung asil
ping-pingan $\alpha$ karo~$\beta$. Kanggo kuwi, bayangna kothak-kothak
persegi panjang kang ambané $\alpha$ lan dhuwuré $\beta$; asil
ping-pingan $\alpha$ lan $\beta$ dipikolehi kanthi ngliwati saben
larik, banjur ngliwati larik sabanjuré\ldots nganti kabèh larik
rampung diliwati. Kanthi tembung liya, asil ping-pingan $\alpha$
lan $\beta$ dipikolehi kanthi ngganti \emph{saben} unsur ing
$\beta$ nganggo salinan $\alpha$.

Kanggo netepaké kanthi formal, kita mung migunakaké pengurutan
leksikografis kuwalik ing ping-pingan Kartesius saka $\alpha$ lan $\beta$:

\begin{defn}
Kanggo sembarang ordinal $\alpha, \beta$, asil ping-pingané yaiku $\alpha \ordtimes \beta = \ordtype{\alpha \times \beta, \rlexless}$.
\end{defn}
\noindent
Kita kudu mesthèkaké manèh yèn definisi iki wujudé trep:

\begin{lem}\ollabel{ordtimeslessiswo}
$\tuple{\alpha \times \beta, \rlexless}$ iku pengurutan becik kanggo
sembarang ordinal $\alpha$ lan $\beta$.
\end{lem}

\begin{proof}
Persis kaya bukti \olref[add]{ordsumlessiswo}.
\end{proof}
\noindent
Ora angel uga mbuktèkaké yèn ping-pingan manut pratelan iki:

\begin{lem}\ollabel{ordtimesrecursion}
Kanggo sembarang ordinal $\alpha, \beta$:
\begin{align*}
	\alpha \ordtimes 0 &= 0\\
	\alpha \ordtimes (\beta \ordplus 1) &=
		(\alpha \ordtimes \beta) \ordplus \alpha\\
	\alpha  \ordtimes \beta &=
		\supstrict_{\delta < \beta}(\alpha \ordtimes \delta) &&
		\text{yèn $\beta$ ordinal limit}.
\end{align*}
\end{lem}

\begin{proof}
Dipasrahaké minangka latihan.
\end{proof}

Pancen, kaya ing kasus panambahan, kita bisa netepaké ping-pingan
ordinal liwat persamaan rekursi iki, tanpa menehi definisi langsung.
Uga kaya panambahan, sawetara sipaté wis kita kenal:

\begin{lem}\ollabel{ordinalmultiplicationisnice}
Yèn $\alpha, \beta, \gamma$ ordinal, mula:
\begin{enumerate}
	\item\ollabel{ordtimes1} yèn $\alpha \neq 0$ lan $\beta < \gamma$,
	mula $\alpha \ordtimes \beta < \alpha \ordtimes \gamma$;
	\item\ollabel{ordtimes2} yèn $\alpha \neq 0$ lan $\alpha \ordtimes
	\beta = \alpha\ordtimes\gamma$, mula $\beta = \gamma$;
	\item\ollabel{ordtimes3} $\alpha \ordtimes (\beta \ordtimes
	\gamma) = (\alpha \ordtimes \beta) \ordtimes \gamma$;
	\item\ollabel{ordtimes4} yèn $\alpha \leq \beta$, mula $\alpha
	\ordtimes \gamma \leq \beta \ordtimes \gamma$;
	\item\ollabel{ordtimes5} $\alpha \ordtimes (\beta \ordplus
	\gamma) = (\alpha \ordtimes \beta )\ordplus (\alpha\ordtimes
	\gamma)$.
\end{enumerate}
\end{lem}

\begin{proof}
Dipasrahaké minangka latihan.
\end{proof}

Panjenengan bisa mbuktèkaké utawa nggoleki asil-asil liya
sakarepé. Nanging, sawisé
\olref[ord-arithmetic][add]{ordsumnotcommute}, pratelan
iki mesthiné ora nggumunaké:

\begin{prop}
Ping-pingan ordinal ora komutatif: $2 \ordtimes \omega  =
\omega < \omega \ordtimes 2$
\end{prop}

\begin{proof}
$2 \ordtimes \omega = \supstrict_{n < \omega}(2\ordtimes  n) = \omega \in \supstrict_{n < \omega}(\omega \ordplus n) = \omega \ordplus \omega = \omega \ordtimes 2$.
\end{proof}
\noindent
Maneh, panjelasan intuitifé prasaja. Kanggo ngitung $2
\ordtimes \omega$, saben wilangan asli diganti karo rong unsur.
Jinis urutané bakal padha yèn kabèh wilangan “setengah” diselipaké
ing antarané wilangan asli, yaiku yèn kita ndeleng urutan alami ing
$\Setabs{\nicefrac{n}{2}}{n \in \omega}$. Kanthi cara ndeleng iki,
jinis urutané cetha padha karo $\omega$ dhéwé. Nanging kanggo
ngitung $\omega \ordtimes 2$, kita nyelehake rong salinan $\omega$,
siji sawisé sijiné.

\begin{prob}
Buktèkna
\olref[sth][ord-arithmetic][mult]{ordtimeslessiswo},
\olref[sth][ord-arithmetic][mult]{ordtimesrecursion},
lan
\olref[sth][ord-arithmetic][mult]{ordinalmultiplicationisnice}
\end{prob}

\end{document}
